\chapter{Current Evidence Boundary}
\index{evidence!current boundary|textbf}

\label{ch:results-and-limitations}

\chapterstatus{retained evidence inventory and result-reporting framework. No
production bulk, impurity, tight-binding, or continuum result is accepted.}

This chapter records what the project can presently support and, equally
importantly, what it cannot. Current status is derived from the owning task,
calculation, and evidence records associated with this draft. Those records
remain authoritative if later work supersedes this snapshot.

\section{Meaning of ``result'' in this chapter}

A result is reported with its evidence class and provenance. The categories used
here are

\begin{itemize}
  \item proved or mechanized mathematical result with declared assumptions and
    proof status;
  \item verified software behavior;
  \item verified numerical behavior on controlled mathematics;
  \item tutorial or bootstrap execution observation;
  \item calculated physical result with retained provenance;
  \item scientific validation result; and
  \item uncertainty-quantification result.
\end{itemize}

These categories are not ordered stages through which an array advances
automatically. A proof establishes only its implication from its assumptions,
and the same artifact may support one bounded software claim while
remaining unusable for a physical claim.

\section{Accepted represented-operator software foundation}
\index{represented operator!software foundation}


\evidenceclass{software verification, with declared controlled
numerical-verification cases.}

The proof registry separately records nine frozen finite-dimensional
\texttt{PRF-05} contracts. Only projector invariance under unitary frame
rotation, \texttt{PRF-05.01}, has been checked by one pinned Lean backend. The
other Lean targets and every Isabelle and Rocq target remain unencoded; no
contract is cross-checked and all physics proof packages remain proposed. These
facts are mechanization-status results, not physical validation.

The accepted Python operator-record foundation provides

\begin{itemize}
  \item finite represented-operator storage with explicit state-space, basis,
    geometry, energy-reference, provenance, and matrix metadata;
  \item Hermiticity analysis in a fixed representation;
  \item deterministic version-one JSON serialization with a public schema and
    golden fixtures;
  \item exact compatibility auditing for represented metadata;
  \item signed subtraction for records that are already compatible;
  \item maximum-entry, Frobenius, and spectral residual analysis; and
  \item a composed comparison workflow for differencing and residual analysis.
\end{itemize}

This is a software capability result, not an accepted scientific operator. The
foundation does not align bases or gauges, convert units, align scalar energy
zeros, transform geometries, establish physical equivalence, or decide that a
generic difference is an impurity operator. It has not supplied scientific
validation or uncertainty quantification, and it does not by itself pass the
later bulk or alignment gates.

\section{Accepted legacy tutorial observations}
\index{evidence!legacy tutorial}


\evidenceclass{tutorial execution, artifact provenance, and bounded
software-extraction behavior.}

Four legacy silicon tutorial tasks are retained as human-accepted bounded
evidence: an SCF tutorial reproduction, artifact inventory, periodic-record
extraction, and a matching tutorial band calculation. The band execution
consumed an isolated identity-verified copy of the accepted tutorial SCF state
and retained 28 ordered $\mathbf k$ points with eight bands at each point.

Against the bundled legacy printed reference, 166 of 224 eigenvalues matched at
printed precision and the largest printed difference was approximately
$0.0001\,\mathrm{eV}$. No numerical comparison tolerance or pass/fail rule was
accepted, so this observation is not promoted into numerical verification. An
IEEE floating-point exception report remains unresolved and unclassified.

The tutorial parent uses a legacy LDA pseudopotential and tutorial numerical
settings. Its energy reference, basis identity, retained subspace, and gauge are
insufficient for the production representation and reduction questions. These
records therefore establish neither a production silicon band structure nor a
band gap, valley position, effective mass, Wannier target, or tight-binding
reference.

\section{Direct PBE bootstrap observations}
\index{evidence!PBE bootstrap}


\evidenceclass{bootstrap scientific-harness development observation.}

A separately retained PBE/PseudoDojo bootstrap campaign executed nine SCF and
nine linked NSCF Quantum ESPRESSO invocations. Every invocation exited with
status zero and emitted \texttt{JOB DONE.}; no retry was performed. Exact input,
executable, and pseudopotential identities and compact finite-setting records
are retained by the owning convergence design and disposition.

\citationtodo{high priority}{Attach the standard scholarly sources for PBE,
PseudoDojo, and Quantum ESPRESSO---keys \texttt{perdew1996},
\texttt{vansetten2018}, and \texttt{giannozzi2009} (already present)---only
after checking their local scope. These sources identify methods and software;
they do not substantiate the repository-specific execution count or outcomes,
which require durable internal provenance.}

These observations demonstrate that the identified cases executed under the
direct bootstrap boundary. They are not canonical
\texttt{ScientificWorkflowRun} evidence. No final wavefunction cutoff,
charge-density cutoff, $\mathbf k$ mesh, lattice parameter, indirect gap,
valley coordinate, or effective mass was selected. No infinite-basis estimate,
production-parent acceptance, scientific validation, or UQ result follows from
the successful exits.

The next convergence execution is deferred pending scientific-harness
reimplementation and separate authorization. The retained observations may
inform that implementation without being discarded or silently reclassified.

\section{Controlled periodic-reduction benchmarks}
\index{evidence!periodic reduction benchmarks}

\evidenceclass{human-accepted calculated illustrative numerical verification.}

Appendix~\ref{app:one-dimensional-reduction} reports a controlled
one-dimensional cosine-lattice reduction. Independently refined plane-wave and
finite-difference representations converge toward the same low-mode parent;
selected energies agree with Mathieu characteristic values; and isolated- and
composite-band constructions verify parallel transport, alignment, hopping
reconstruction, and finite-range truncation. A separate Wannier90 calculation
uses the same eigenvalue and overlap interface and reproduces the aligned
represented operators after convergence.

The benchmark also identifies bounded failure modes. Scalar energy-ordered
labels become poorly conditioned near adjacent-gap closure, rough gauges degrade
finite-range locality without changing the complete spectrum, and direct and
mediated reductions cease to agree when their weights or fitting domains differ.
These are numerical-verification results for the synthetic cosine parent. They
are not a production silicon calculation, scientific validation, uncertainty
quantification, or evidence of transferability to silicon. At the evidence
boundary represented by this draft, the calculated records pass their declared
verification checks and the bounded one-dimensional periodic-reduction task has
received final human acceptance under these stated limitations.

Appendix~\ref{app:two-dimensional-reduction} extends the controlled periodic
program to a two-dimensional synthetic parent. The accepted evidence covers
separable and coupled scalar operators, a corrected composite-band polar gauge,
mass-tensor recovery, real-space hopping reconstruction, common-estimator
localization comparisons, distinct QWZ, Hofstadter, and Haldane topology
controls, and bounded mesh, cutoff, embedding, and hopping sensitivities.
Native Berry-link spreads and common finite-supercell spread estimates remain
separate quantities.

The expanded optimizer study retains 256 frozen start/interface trajectories.
Their final effective endpoints comprise 196 native convergences and 60
iteration-limit nonconvergences. The frozen mesh-convergence and repeated-basin
criteria are not supported, and the density-aware symmetry analysis leaves
every frozen basin as a singleton. The right-censored convergence-time
regression is exploratory and descriptive. Only one trial-projection family
was executed; projection-family sensitivity is explicitly excluded from the
closed benchmark.

This is human-accepted negative numerical evidence for one frozen synthetic
protocol. It does not establish global optimization, general Wannier
convergence, uncertainty quantification, a DFT interface, or material
validation.

\section{Controlled impurity-extraction benchmarks}
\index{evidence!impurity extraction benchmark}

\evidenceclass{human-accepted calculated synthetic software and numerical
verification, plus a human-accepted bounded literature synthesis.}

Appendix~\ref{app:impurity-model-classes} reports a spin-space diagnostic, a
matched one-dimensional pristine--defect extraction, and four follow-ons that
separate blind alignment, independent extraction routes, a finite-rank
resolvent oracle, and continuum refinement. The numerical records verify their
declared finite-model contracts and preserve failed compatibility controls,
finite-volume effects, operator residuals, spectral errors, projector errors,
and continuum axes as separate quantities. They do not supply a production
silicon operator or a validated material model.

The accompanying primary-literature comparison narrows the contribution
boundary. Finite-rank impurity methods, weak localization, band-edge
homogenization, explicit discrete--continuum identifications, periodic-volume
effects, multivalley and central-cell donor models, Wannier defect operators,
and gauge-aligned model stitching all have established or contemporary
precedents~\cite{kosterSlater1954impurity,simon1976weak,parzygnat2010,hoeferWeinstein2011,ducheneVukicevicWeinstein2015,nakamuraTadano2021,koenigLeeHammer2011,kohnLuttinger1955donor,gamble2015,berlijnVoljaKu2011,corsettiMostofi2011,lihmPark2019,luParkZhouBernardi2020}.
Two relevant 2026 records are explicitly preprints: an onsite-only
substitutional-defect screening protocol and a clean/defect Wannier workflow
using assignment, unitary Procrustes alignment, and a nontrivial overlap
metric~\cite{kangMuechler2026,shiXieZhangChuGong2026}.

The bounded contribution is therefore a reproducible integration of these
concerns into one fail-closed synthetic diagnostic chain, not a new component
theorem, a priority claim, or scientific validation. The review's negative
search outcome is conditional on its declared concepts, queries, interfaces,
dates, and stopping rule; it is not evidence that no unreviewed precedent
exists.

Appendix~\ref{app:two-dimensional-defect-extraction} separately records the
two-dimensional defect and material-transfer sequence. The controlled synthetic
defect task is active and consumes the accepted periodic parent and spin-space
and one-dimensional extraction contracts, but it contains no completed
two-dimensional defect calculation. Material transfer remains proposed and
inactive.

\section{Current scientific result boundary}
\index{result!scientific boundary}


At the status represented by this draft, the project has not accepted

\begin{itemize}
  \item a converged production bulk-silicon Kohn--Sham parent;
  \item a production equilibrium lattice parameter;
  \item a calculated production indirect gap or conduction-valley position;
  \item longitudinal or transverse production electron effective masses;
  \item a validated bulk-silicon Wannier operator;
  \item a fitted or selected $sp^3s^{\ast}$ tight-binding model;
  \item a frozen basis- and gauge-alignment protocol;
  \item a calculated two-dimensional defect-extraction or material-transfer
    result;
  \item a phosphorus or boron first-principles impurity operator;
  \item a minimal impurity lattice model;
  \item a donor or acceptor continuum crossover radius;
  \item scientific validation for a declared physical use; or
  \item propagated uncertainty for the proposed reduction hierarchy.
\end{itemize}

Accordingly, this chapter contains no production band plot, fitted parameter
table, impurity profile, binding energy, or crossover curve. Adding an expected
shape or literature value to fill visual space would obscure the evidentiary
boundary.

\section{Present limitations}
\index{limitation!current}


\subsection{Parent calculation}

The physical and numerical protocols are accepted, but the production parent
has not crossed its convergence and acceptance boundary. Finite candidate
settings and successful process completion cannot substitute for that gate.
The selected PBE parent will also retain parent-model limitations relative to
experiment even after numerical convergence.

\citationtodo{medium priority}{Narrow this broad statement to a named observable
before citing it. If the intended claim is the electronic gap, check
Perdew--Levy (1983) and Sham--Schluter (1983), bibliography keys
\texttt{perdewLevy1983} and \texttt{shamSchluter1983} (already present). Do not
use those sources to imply a universal PBE limitation.}

\subsection{Representation}

For production silicon, the target Wannier rank, projections, energy windows,
uniform interface grid, localization outcome, and interpolation domain have not
been accepted. The accepted synthetic benchmarks do not select or validate
those material settings. The alignment map between independently generated
spaces is likewise unavailable. Operator subtraction beyond already compatible
synthetic or software records would therefore lack required meaning.

\subsection{Bulk model reduction}

The executable $sp^3s^{\ast}$ hierarchy, frozen training and withheld sets,
alignment group, shell weights, objective normalizations, operator tolerance,
and complete model-selection rule remain downstream work. Proposed manuscript
tolerances are not treated as accepted computational evidence merely because
they appear in a draft.

\subsection{Impurity branches}

Phosphorus and boron specialization gates have not supplied their production
pseudopotentials, supercell sequences, doped Wannier representations, or
alignment records. Periodic-image convergence is unavailable. The boron final
branch additionally requires compatible SOC and spinor treatment, while the
phosphorus neutral branch requires its open-shell and charge-state
interpretation to remain explicit.

\subsection{Continuum reduction}

Verified continuum solver infrastructure would not by itself establish a
first-principles reduction. Host parameters, accepted atomistic impurity
operators, continuum-to-Wannier embedding, spatial error criteria, and
bound-state validation are all required before a crossover statement can be
made. Established finite-rank identities, band-edge homogenization results,
discrete-to-continuum convergence theorems, and semiconductor effective-mass
models provide methodological precedents, but citations to those results do not
establish that this project's represented impurity operator satisfies their
hypotheses or has reached a material continuum crossover.

\section{Negative findings}
\index{negative finding}


The failed periodic-2D localization criteria are accepted negative numerical
evidence for their frozen synthetic protocol. They are not a scientific
negative finding about silicon, defects, or the proposed material model
hierarchy.

No scientific negative finding has yet been accepted for the model hierarchy.
In particular, the project has not shown that

\begin{itemize}
  \item nearest-neighbor $sp^3s^{\ast}$ is incompatible with the joint criteria;
  \item a larger tight-binding class is required;
  \item no reduced impurity model passes;
  \item no finite continuum crossover exists; or
  \item two declared reduction paths fail to commute beyond tolerance.
\end{itemize}

These are possible future outcomes, not conclusions. When a frozen tested domain
produces such an outcome, the result should state the domain, metric, tolerance,
optimization or convergence evidence, and unresolved alternatives.

\section{Result table contract}
\index{result table!contract}


Each future numerical table should identify

\begin{itemize}
  \item the calculated quantity and units;
  \item immutable parent, child, and analysis identities;
  \item representation, basis, gauge, geometry, spin, and energy reference where
    applicable;
  \item numerical settings and tested refinement domain;
  \item uncertainty or observed sensitivity without false precision;
  \item comparison reference and evidence class;
  \item metric, tolerance, and pass/fail rule; and
  \item limitations on interpretation.
\end{itemize}

Training results and withheld results appear in separate columns or tables.
Parent-model, numerical, and model-reduction errors remain separate quantities.
A missing entry is marked unavailable rather than estimated from an unrelated
source.

\section{Figure contract}
\index{figure!contract}


A future figure must be a deterministic view of retained data. Its caption and
associated record identify the data source, processing code, axes, units,
reference zero, plotted subset, and evidentiary status. Expected diagrams,
synthetic regressions, tutorial bands, and production results use visibly
distinct labels.

Band plots state the reciprocal-coordinate path and energy reference. Operator
plots state the basis, alignment, block or shell definition, norm, and
normalization. Convergence plots show rejected settings and guard points rather
than only the selected candidate. Residual plots distinguish training from
withheld data.

\section{Path to populated scientific results}

The first populated production sequence requires an accepted bulk convergence
record, equilibrium lattice reference, and production SCF parent. Purpose-built
band-path and local-valley children then provide the bulk observables. A
separately approved uniform NSCF and QE--Wannier interface support candidate
Wannier operators, followed by interpolation and alignment evidence. Tight-
binding, impurity, and continuum result sections become eligible only after
their corresponding parent artifacts exist.

Until those records are available, the absence of production numbers in this
chapter is an accurate result boundary rather than an incomplete editorial
placeholder.
